\documentclass[11pt, a4paper]{article}

\usepackage[T1]{fontenc}
\usepackage[utf8]{inputenc}
\usepackage{tgpagella}
\usepackage[spanish, es-noshorthands]{babel}
\usepackage{amsmath, amssymb, bm}
\usepackage{tabularx}
\usepackage{booktabs}
\usepackage[most]{tcolorbox}
\usepackage[margin=2.5cm]{geometry}
\usepackage{ragged2e}
\usepackage{fancyhdr}

\pagestyle{fancy}
\fancyhf{}
\setlength{\headheight}{16pt}
\lhead{\textbf{UCB Sede Tarija} | Ing. Mecatrónica}
\chead{}
\rhead{IMT-121 Circuitos Electrónicos I | Protocolo Microdefensa}
\lfoot{Prof: M.Sc. Ing. Bernardo Quiroga Turdera}
\rfoot{Página \thepage}
\renewcommand{\headrulewidth}{0.4pt}

\definecolor{ucbblue}{RGB}{0, 51, 102}

\begin{document}

\begin{tcolorbox}[colback=ucbblue!5!white, colframe=ucbblue, title=\textbf{Guía de Ejecución: Auditoría por Microdefensa}]
    \centering \textbf{Uso Exclusivo del Instructor}
\end{tcolorbox}

\section*{1. Naturaleza y Propósito}
La microdefensa \textbf{no es un tercer componente de evaluación} (carece de peso fijo independiente), no es una evaluación universal ni un castigo por reprobar. Es una herramienta de auditoría oral (duración acotada) para resolver \textbf{discrepancias documentadas} entre la evidencia entregada por el estudiante.

\section*{2. Criterios para Activar el Trigger}
Se activa la microdefensa SÓLO si existe una anomalía en la "Matriz de Comparación Post-Evaluación", por ejemplo:
\begin{itemize}
    \item \textbf{Disparidad Cognitiva Extrema:} El reporte IEEE del equipo presenta un análisis impecable y complejo del "Efecto de Carga", pero en el examen escrito individual el estudiante no logró identificar ni el puerto de salida.
    \item \textbf{Asimetría de Autoría:} Se sospecha fuertemente que uno de los integrantes del equipo no participó en la construcción ni análisis de los datos del laboratorio.
    \item \textbf{Incompatibilidad Práctica:} El reporte incluye una curva de potencia obtenida con mediciones físicas de laboratorio, pero la evidencia práctica en sala indicaba que el equipo jamás logró aislar el cortocircuito.
\end{itemize}

\section*{3. Protocolo de Ejecución}
\begin{enumerate}
    \item Identifique la incertidumbre exacta.
    \item Formule entre 2 y 4 preguntas técnicas cerradas, atadas a la evidencia en duda (Ej. \textit{"En el reporte entregado afirman haber medido $R_{th}$ con la red desenergizada. Explíqueme por qué apagar el banco es indispensable para obtener ese valor y no otro"}).
    \item Registre el resultado de la auditoría: \textbf{Confirmada} (dominio demostrado), \textbf{Clarificada} (malentendido aislado), o \textbf{No Resuelta} (evidencia del reporte rechazada para ese individuo).
    \item Cualquier afectación a la calificación se aplica sobre el porcentaje del "Componente de Laboratorio", ajustando el peso de acuerdo con la invalidez de la evidencia rechazada.
\end{enumerate}

\end{document}
